\section{AIME I 1983}
        \begin{enumerate}
        	\item (\textit{AIME I 1983}) Let $x$, $y$ and $z$ all exceed $1$ and let $w$ be a positive number such that $\log_xw=24$, $\log_y w = 40$ and $\log_{xyz}w=12$. Find $\log_zw$.


 

	\item (\textit{AIME I 1983}) Let $f(x)=|x-p|+|x-15|+|x-p-15|$, where $0 < p < 15$. Determine the \href{https://artofproblemsolving.com/wiki/index.php?title=Minimum}{minimum} value taken by $f(x)$ for $x$ in the \href{https://artofproblemsolving.com/wiki/index.php?title=Interval}{interval} $p \leq x\leq15$.


 

	\item (\textit{AIME I 1983}) What is the product of the real roots of the equation $x^2 + 18x + 30 = 2 \sqrt{x^2 + 18x + 45}$?


 

	\item (\textit{AIME I 1983}) A machine-shop cutting tool has the shape of a notched circle, as shown. The radius of the circle is $\sqrt{50}$ cm, the length of $AB$ is $6$ cm and that of $BC$ is $2$ cm. The angle $ABC$ is a right angle. Find the square of the distance (in centimeters) from $B$ to the center of the circle.


 
                    \begin{figure}[H]
                        \centering
                        \includegraphics[width=0.5\textwidth]{Problem/AIME_I_1983_4_Problem_Figure_1.png}
                    \end{figure}
                    


 

	\item (\textit{AIME I 1983}) Suppose that the sum of the squares of two complex numbers $x$ and $y$ is $7$ and the sum of the cubes is $10$. What is the largest real value that $x + y$ can have?


 ayyy 67


 

	\item (\textit{AIME I 1983}) Twenty five of King Arthur's knights are seated at their customary round table. Three of them are chosen - all choices being equally likely - and are sent off to slay a troublesome dragon. Let $P$ be the probability that at least two of the three had been sitting next to each other. If $P$ is written as a fraction in lowest terms, what is the sum of the numerator and denominator?


 

	\item (\textit{AIME I 1983}) What is the smallest $2$-digit prime factor of the integer $n = {200\choose 100}$?


 

	\item (\textit{AIME I 1983}) Find the minimum value of $\frac{9x^2\sin^2 x + 4}{x\sin x}$ for $0 < x < \pi$.


 

	\item (\textit{AIME I 1983}) The numbers $1447$, $1005$ and $1231$ have something in common: each is a $4$-digit number beginning with $1$ that has exactly two identical digits. How many such numbers are there?


 

	\item (\textit{AIME I 1983}) The solid shown has a square base of side length $s$. The upper edge is parallel to the base and has length $2s$. All other edges have length $s$. Given that $s=6\sqrt{2}$, what is the volume of the solid?


 
                    \begin{figure}[H]
                        \centering
                        \includegraphics[width=0.5\textwidth]{Problem/AIME_I_1983_11_Problem_Figure_1.png}
                    \end{figure}
                    


 

	\item (\textit{AIME I 1983}) Diameter $AB$ of a circle has length a $2$-digit integer (base ten). Reversing the digits gives the length of the perpendicular chord $CD$. The distance from their intersection point $H$ to the center $O$ is a positive rational number. Determine the length of $AB$.


 
                    \begin{figure}[H]
                        \centering
                        \includegraphics[width=0.5\textwidth]{Problem/AIME_I_1983_12_Problem_Figure_1.png}
                    \end{figure}
                    


 

	\item (\textit{AIME I 1983}) For $\{1, 2, 3, \ldots, n\}$ and each of its nonempty subsets a unique \textbf{alternating sum} is defined as follows. Arrange the numbers in the subset in decreasing order and then, beginning with the largest, alternately add and subtract succesive numbers. For example, the alternating sum for $\{1, 2, 3, 6,9\}$ is $9-6+3-2+1=5$ and for $\{5\}$ it is simply $5$. Find the sum of all such alternating sums for $n=7$.


 

	\item (\textit{AIME I 1983}) In the adjoining figure, two circles with radii $8$ and $6$ are drawn with their centers $12$ units apart. At $P$, one of the points of intersection, a line is drawn in such a way that the chords $QP$ and $PR$ have equal length. Find the square of the length of $QP$. 


 
                    \begin{figure}[H]
                        \centering
                        \includegraphics[width=0.5\textwidth]{Problem/AIME_I_1983_14_Problem_Figure_1.png}
                    \end{figure}
                    


 

	\item (\textit{AIME I 1983}) The adjoining figure shows two intersecting chords in a circle, with $B$ on minor arc $AD$. Suppose that the radius of the circle is $5$, that $BC=6$, and that $AD$ is bisected by $BC$. Suppose further that $AD$ is the only chord starting at $A$ which is bisected by $BC$. It follows that the sine of the central angle of minor arc $AB$ is a rational number. If this number is expressed as a fraction $\frac{m}{n}$ in lowest terms, what is the product $mn$?

                    \begin{figure}[H]
                        \centering
                        \includegraphics[width=0.5\textwidth]{Problem/AIME_I_1983_15_Problem_Figure_1.png}
                    \end{figure}
                    


 

\end{enumerate}